\chapter*{Appendix O: Semantic Logic}
\addcontentsline{toc}{chapter}{Appendix O: Semantic Logic}

\section{Syntax of Semantic Propositions}

\subsection{Semantic Signatures}

Let $\Sigma$ be a signature consisting of:
\begin{itemize}
\item a set of types $\mathsf{Ty}$,
\item a set of operation symbols $\mathsf{Op}$,
\item a set of relation symbols $\mathsf{Rel}$ with arities in $\mathsf{Ty}^\ast$.
\end{itemize}

\subsection{Formulas}

The set of well-formed formulas $\mathrm{Form}(\Sigma)$ is generated by:
\[
\varphi ::= R(t_1,\dots,t_n) \mid \varphi\land\varphi \mid \varphi\lor\varphi \mid \varphi\to\varphi \mid \lnot \varphi \mid \forall x:\tau.\,\varphi \mid \exists x:\tau.\,\varphi,
\]
where $R\in\mathsf{Rel}$ and $t_i$ are terms over $\Sigma$.

\section{Semantic Structures}

\subsection{Interpretations}

A semantic structure for $\Sigma$ is a tuple
\[
\mathcal{M} = (M_\tau, \, O_f, \, R_g)
\]
where:
\begin{itemize}
\item for each $\tau\in\mathsf{Ty}$, $M_\tau$ is a semantic domain,
\item for each $f\in\mathsf{Op}$ with arity $\tau_1\times\cdots\times\tau_n\to\tau$,  
  $O_f : M_{\tau_1}\times\cdots\times M_{\tau_n}\to M_\tau$,
\item for each $g\in\mathsf{Rel}$ with arity $\tau_1\times\cdots\times\tau_n$,  
  $R_g \subseteq M_{\tau_1}\times\cdots\times M_{\tau_n}$.
\end{itemize}

\subsection{Assignments}

An assignment $\rho$ maps each variable $x:\tau$ to an element $\rho(x)\in M_\tau$.

\section{Satisfaction Relation}

\subsection{Definition}

The satisfaction relation is defined inductively:
\begin{align*}
\mathcal{M},\rho \models R(t_1,\dots,t_n)
&\quad\Longleftrightarrow\quad
(\llbracket t_1 \rrbracket_\rho,\dots,\llbracket t_n \rrbracket_\rho) \in R,\\[4pt]
\mathcal{M},\rho \models \varphi\land\psi
&\quad\Longleftrightarrow\quad
\mathcal{M},\rho\models\varphi\ \text{and}\ \mathcal{M},\rho\models\psi,\\[4pt]
\mathcal{M},\rho \models \varphi\lor\psi
&\quad\Longleftrightarrow\quad
\mathcal{M},\rho\models\varphi\ \text{or}\ \mathcal{M},\rho\models\psi,\\[4pt]
\mathcal{M},\rho \models \varphi\to\psi
&\quad\Longleftrightarrow\quad
\text{if }\mathcal{M},\rho\models\varphi\text{ then }\mathcal{M},\rho\models\psi,\\[4pt]
\mathcal{M},\rho \models \lnot\varphi
&\quad\Longleftrightarrow\quad
\mathcal{M},\rho\nvDash\varphi,\\[4pt]
\mathcal{M},\rho \models \forall x\!:\!\tau.\,\varphi
&\quad\Longleftrightarrow\quad
\forall a\in M_\tau.\ \mathcal{M},\rho[x\mapsto a]\models\varphi,\\[4pt]
\mathcal{M},\rho \models \exists x\!:\!\tau.\,\varphi
&\quad\Longleftrightarrow\quad
\exists a\in M_\tau.\ \mathcal{M},\rho[x\mapsto a]\models\varphi.
\end{align*}

\section{Semantic Consequence}

\subsection{Definition}

For a set of formulas $\Gamma$ and formula $\varphi$,
\[
\Gamma \models \varphi
\quad\Longleftrightarrow\quad
\text{for all }\mathcal{M},\rho,\ 
\mathcal{M},\rho\models\Gamma\Rightarrow\mathcal{M},\rho\models\varphi.
\]

\subsection{Equivalence}

Two formulas are semantically equivalent if:
\[
\varphi \equiv \psi
\quad\Longleftrightarrow\quad
\forall\mathcal{M},\rho.\ 
(\mathcal{M},\rho\models\varphi \ \text{iff}\ \mathcal{M},\rho\models\psi).
\]

\section{Modal Operators for Semantic Flow}

\subsection{Flow Modality}

Let $F:\mathbf{RSVP}\to\mathbf{RSVP}$ be a semantic flow endofunctor.

Define modal operators:
\[
\Box\varphi \quad\text{and}\quad \Diamond\varphi
\]
by:
\[
\mathcal{M},\rho\models\Box\varphi
\quad\Longleftrightarrow\quad
F(\mathcal{M}),\rho\models\varphi,
\]
\[
\mathcal{M},\rho\models\Diamond\varphi
\quad\Longleftrightarrow\quad
\exists \mathcal{N},\iota.\  
\iota:\mathcal{M}\to F(\mathcal{N})\ \text{and}\ \mathcal{N},\rho\models\varphi.
\]

\section{Fixed-Point Logic}

\subsection{Least and Greatest Fixed Points}

Let $\Phi$ be a monotone operator on formulas:
\[
\Phi : \mathrm{Form}\to\mathrm{Form}.
\]

Define:
\begin{align*}
\mu X.\,\Phi(X) &= \text{least fixed point of }\Phi,\\
\nu X.\,\Phi(X) &= \text{greatest fixed point of }\Phi.
\end{align*}

\subsection{Satisfaction}

\[
\mathcal{M},\rho\models\mu X.\,\Phi(X)
\quad\Longleftrightarrow\quad
\mathcal{M},\rho\models\Phi(\mu X.\,\Phi(X)),
\]
\[
\mathcal{M},\rho\models\nu X.\,\Phi(X)
\quad\Longleftrightarrow\quad
\mathcal{M},\rho\models\Phi(\nu X.\,\Phi(X)).
\]

\section{Inference System}

\subsection{Axioms}

\begin{itemize}
\item All propositional tautologies.
\item Quantifier axioms for $\forall$ and $\exists$.
\item Modal distribution axioms for $\Box$ and $\Diamond$.
\item Fixed-point unfolding axioms:
\[
\mu X.\,\Phi(X) \leftrightarrow \Phi(\mu X.\,\Phi(X)),
\]
\[
\nu X.\,\Phi(X) \leftrightarrow \Phi(\nu X.\,\Phi(X)).
\]
\end{itemize}

\subsection{Rules}

\begin{itemize}
\item Modus ponens.
\item Generalisation:  
  $\varphi \vdash \forall x.\,\varphi$.
\item Monotone induction for $\mu$:
  \[
  \Phi(\varphi) \vdash \varphi
  \ \Longrightarrow\
  \mu X.\,\Phi(X) \vdash \varphi.
  \]
\item Coinduction for $\nu$:
  \[
  \varphi \vdash \Phi(\varphi)
  \ \Longrightarrow\
  \varphi \vdash \nu X.\,\Phi(X).
  \]
\end{itemize}

\section{Completeness}

\subsection{Semantic Completeness}

For the above inference system,
\[
\Gamma \vdash \varphi
\quad\Longleftrightarrow\quad
\Gamma \models \varphi
\]
holds for $\Sigma$-structures in $\mathbf{RSVP}$.

\subsection{Conservativity}

If $\varphi$ is purely first-order (no modality, no fixed points),  
the semantic logic reduces to standard first-order logic with:
\[
\Gamma\models_{\mathbf{RSVP}}\varphi
\iff
\Gamma\models_{\mathrm{FOL}}\varphi.
\]

